% Add "answers" option to print the answers
%\documentclass[addpoints,answers]{exam}
\documentclass[addpoints]{exam}

\usepackage[utf8]{inputenc}
\usepackage{calc}
\usepackage{listings}
\usepackage{xcolor}
\usepackage{xspace}

\definecolor{codegreen}{rgb}{0,0.6,0}
\definecolor{codegray}{rgb}{0.5,0.5,0.5}
\definecolor{codepurple}{rgb}{0.58,0,0.82}

\lstdefinestyle{Cstyle}{
    commentstyle=\color{codegreen},
    keywordstyle=\color{magenta},
    numberstyle=\tiny\color{codegray},
    stringstyle=\color{codepurple},
    basicstyle=\ttfamily\small,
    breakatwhitespace=true,         
    breaklines=true,                 
    captionpos=b,                    
    keepspaces=true,                 
    numbers=left,                    
    numbersep=5pt,                  
    showspaces=false,                
    showstringspaces=false,
    showtabs=false,                  
    tabsize=4
}

\lstset{style=Cstyle}

% \makebox[\textwidth]{Name:\enspace\hrulefill}
%  
% \vspace{5mm}
%\makebox[\textwidth]{Instructor’s name:\enspace\hrulefill}

\newcommand{\truefalse}
{		\begin{checkboxes}
			 \choice True
			 \choice False. Explanation: 
		\end{checkboxes}
		\vspace{10pt}
}

\newcommand{\answerspace}[1]
{
	\ifprintanswers
	\else
		\setlength{\mylength}{10pt}%
		%\noindent\rule{\mylength}{20pt}
		%\vspace{\fpeval{1.2*(#1)}}
		\vspace{#1\mylength}
	\fi
}
\newlength{\mylength}

\newcommand\pf[1]{\textcolor{red}{PF: #1}}
\newcommand\am[1]{\textcolor{blue}{AM: #1}}

\newcommand{\answer}[1]{\begin{solution}{#1}\end{solution}}
\newcommand{\ra}{$\rightarrow$\xspace}



%\usepackage{listings}
\usepackage{color}
\usepackage{lastpage}
\usepackage{verbatim}
\usepackage{fancyvrb}
\usepackage{graphicx}
\usepackage{float}
\usepackage{algorithm}
\usepackage{algpseudocode}
\usepackage{minted}

%\pagestyle{fancy}
%\fancyfoot[C]{{\thepage} of \pageref{LastPage}}
\newcommand{\includecode}[2][c]{\lstinputlisting[caption=, escapechar=, language=#1]{#2}<!---->}
\RecustomVerbatimCommand{\VerbatimInput}{VerbatimInput}{}

\renewcommand{\solutiontitle}{\noindent\textbf{Solution sketch: }}
\renewcommand{\pf}[1]{}

\newcommand{\boxedanswer}[2]{\ifprintanswers\answer{#2}\else\makeemptybox{#1}\fi}

\definecolor{SolutionColor}{rgb}{0.8,0.9,1}


\cfoot{Page \thepage\ of \pageref{LastPage}}

\rhead{Initials: \rule{20mm}{0.4pt}}



%%%%%%%%%%%%%%%%%%%%%%%%%%%%%%%%%%%%%%%%
\begin{document}
\fvset{numbers=left,numbersep=3pt}
\thispagestyle{empty}

\begin{center}
  {\LARGE CS408 Fall 2026 -- Homework 1\\}

  {\Large Deadline: Friday, September 11th \textbf{at 6p} \\}

  \vspace{20mm}
  \fbox{
    \fbox{\parbox{5.5in}{\centering
        \begin{itemize}
        \item The homework is open notes, open book.
        \item No AI tools allowed.
        \item No collaboration and no communication of any kind allowed.
        \item Provide complete and clear but concise answers. Justify thoroughly statements where required.
	\end{itemize}
      }}}
\end{center}

\vspace{13mm}
\begin{center}
	\makebox[0.6\textwidth]{
		\combinedgradetable[h][questions]
	}
\end{center}
%\begin{center}\combinedgradetable[h][questions]\end{center}
\vspace{28mm}

\begin{center}
\makebox[\textwidth]{\textbf{Name (please print)}: \rule{100mm}{0.8pt} }
\end{center}

\newpage
\begin{questions}

\question Consider this faulty program. %

  \begin{minted}{python}
# A function that returns true if the input integer is a positive
# even number, and false otherwise.
def is_positive_even(x: int) -> bool:
    if x < 0:
        return False
    mod_2 = x % 3
    if mod_2 == 0:
        return True
    return False
\end{minted}

\begin{parts}
\part[3] Identify a test case that produces a failure, providing the test
  input, expected behavior, and observed behavior of the program on
  that input. The expected behavior should be consistent with the
  comment at the top of the program.

\boxedanswer{1.5in}{
  \begin{itemize}
  \item \textbf{Input:} \lstinline{x} $\mapsto$ 3
  \item \textbf{Expected output}:  \lstinline{ret} $\mapsto$ false
  \item \textbf{Observed output}:  \lstinline{ret} $\mapsto$ true
  \end{itemize}
}

\part[3] Identify the fault, and fix the fault. %
  \boxedanswer{2in}{ %
    The fault is in line 6: \texttt{mod\_2 = x\%3}. For positive even
    numbers, the check is done by taking modulus of the number with 3,
    which is a wrong way to find if a number is even.  Fix:
    \texttt{mod\_2 = x\%2} }
\part[3] If possible, identify a test case that does not execute the
  fault. %
  \boxedanswer{1.75in}{ %
    Test case: x = -1. If $x<0$, the function does not execute the
    faulty line.  }

  \newpage

\part[3] If possible, identify a test case that executes the fault,
  but whose test path does not include an error state. %
  \boxedanswer{2in}{ %
    Test case: x = 6. If $x = 6$, the fault line executes but the program state remains correct.
  }

\part[3] In the case that an error state is triggered, identify when
  this occurs and describe why that is the first error state.
  \boxedanswer{2in}{ %
    The infection is triggered at line 6:
    \mintinline{python}{mod_2=x\%3}; by definition, the state
    immediately after this fault is executed and infection occurs is
    the first error state.  }

\end{parts}

\newpage

\question Answer the following questions about this pair of faulty and correct programs.

Faulty Code (faulty lines are commented): %
\begin{minted}{python}
# Divide two integers a and b and return the result.
# If b is 0, return -1 to indicate an error.

def divide(a: int, b: int) -> int:
    if b != 0:         # Fault 1
        return 0      # Fault 2
    else:
        return a // b
\end{minted}

Repaired Code: %
\begin{minted}{python}
# Divide two integers a and b and return the result.
# If b is 0, return -1 to indicate an error.

def divide(a: int, b: int) -> int:
    if b == 0:
        return -1
    else:
        return a // b
\end{minted}

\begin{parts}

\part[4] For each faulty line, identify at least one test input that
  makes the line reachable. Explain your reasoning.  %
  \boxedanswer{2in}{%
    The first fault is reachable by \textit{every} test case,
    including, e.g., a test input of a$\mapsto$0 and b $\mapsto$ 2.
      This test input also reaches the second fault.
}

\part [3] For those inputs, explain whether the program state becomes
  infected. Explain your reasoning. %
  \boxedanswer{2in}{%
Indeed, the program state becomes infected for this input, since
  the wrong branch is taken, which manifests as a wrong value for the
  program counter.
}

\part [3] For the inputs you identified, explain whether the infection
  is propagated to the program's output. %
  \boxedanswer{2in}{%
    For this input, the error does not propagate to the
    program's output, since \mintinline{python}{divide} will return
    0, as expected. Note that we got unlucky here, as any other
    value of \mintinline{python}{a} would have lead to a failure.
  }

\end{parts}

\question Consider the python class and one of its tests. You can
  assume any missing code, e.g., \mintinline{python}{Cart},
  \mintinline{python}{Receipt}, behaves as its name suggests.
  \begin{minted}{python}
class PaymentGateway(Protocol):
    # The real gateway honors this contract:
    #   - success:  returns Charge(ok=True)
    #   - declined: RAISES PaymentDeclined  (it never returns ok=False)
    def charge(self, card: Card, amount_cents: int) -> Charge: ...

class Checkout:
    def __init__(self, gateway: PaymentGateway, inventory: Inventory):
        self.gateway = gateway
        self.inventory = inventory
    def purchase(self, cart: Cart) -> Receipt:
        if not self.inventory.reserve(cart):
            return Receipt(status="OUT_OF_STOCK")
        charge = self.gateway.charge(cart.card, cart.total_cents)
        if charge.ok:
            return Receipt(status="CONFIRMED")
        return Receipt(status="DECLINED")

class StubGateway:
    def __init__(self, approves: bool):
        self._approves = approves
    def charge(self, card: Card, amount_cents: int) -> Charge:
        return Charge(id="stub", ok=self._approves)

def test_declined_is_reported() -> None:
    checkout = Checkout(StubGateway(approves=False), InMemoryInventory())
    receipt  = checkout.purchase(Cart(card=BAD_CARD, total_cents=1500))
    assert receipt.status == "DECLINED"
  \end{minted}
\newpage
\begin{parts}
\part[2] Identify the system under test above, and give the line
  numbers for the Arrange / Act / Assert parts of the test. %
  \boxedanswer{1.75in}{
    The system under test is
    \mintinline{python}{checkout.purchase}. Line 26 is the Arrange
    part of the test, line 27 is the Act portion, and line 28 is the
    Assert part.
  }

\part[2] What are the test doubles, and what are regular input values
  in this test? %
  \boxedanswer{1.75in}{
    Here, the test doubles are \mintinline{python}{InMemoryInventory}
    and \mintinline{python}{StubGateway}, and the input value is
    \mintinline{python}{Cart(card=BAD_CARD, total_cents=1500)} }

\part[2] Describe the assumption that \mintinline{python}{StubGateway}
  encodes about the gateway that is inconsistent with the contract for
  \mintinline{python}{PaymentGateway} described in the code comment
  above. %
  \boxedanswer{3.5in}{ %
    \mintinline{python}{StubGateway} assumes a declined card is
    reported by \emph{returning} \mintinline{python}{Charge(ok=False)}
    and that \mintinline{python}{charge} never raises an
    exception. The \mintinline{python}{PaymentGateway} contract says
    the opposite: a decline is reported by \emph{raising}
    \mintinline{python}{PaymentDeclined}, and it \emph{never} returns
    \mintinline{python}{ok=False}. The stub therefore encodes a wrong
    assumption about how a declined payment is signaled.  }

  \newpage
\part[4] Would these inconsistencies lead to a failure when the code
  under test is deployed against \mintinline{python}{PaymentGateway}?
  If so, give the smallest change to the stub that would make
  \mintinline{python}{test_declined_is_reported} catch the fault.
  \boxedanswer{6in}{ %
    Yes. Against \mintinline{python}{PaymentGateway}, a declined card
    makes \mintinline{python}{self.gateway.charge(...)}  \emph{raise}
    \mintinline{python}{PaymentDeclined}. \mintinline{python}{purchase}
    has no handler for it, so the exception propagates out of
    \mintinline{python}{purchase} instead of returning
    \mintinline{python}{Receipt(status="DECLINED")}, and declined
    purchases will fail in production. The unit test misses this
    because the stub returns \mintinline{python}{ok=False} instead of
    raising an exception, so the test never exercises the exception
    path.

    The smallest fix is to simply make the stub honor the contract, by
    raising \mintinline{python}{PaymentDeclined} when
    \mintinline{python}{approves} is \mintinline{python}{False}
    instead of returning \mintinline{python}{Charge(ok=False)}:
% \begin{Verbatim}
% def charge(self, card: Card, amount_cents: int) -> Charge:
%     if not self._approves:
%         raise PaymentDeclined()
%     return Charge(id="stub", ok=True)
% \end{Verbatim}
  With this stub, \mintinline{python}{test_declined_is_reported}
  manifests the exception, \mintinline{python}{purchase} lets the
  exception escape, and the test fails, signalling that
  \mintinline{python}{Checkout} is missing a
  \mintinline{python}{try/except}.  }

\end{parts}

\newpage

\question[5] A tester defined two characteristics for the input
  parameter \texttt{car}, which they used to define the partitioning
  below. Unfortunately, it contains (at least) two mistakes that
  violate the two key properties needed by a valid
  partitioning. Identify the mistake in each block set, name the
  property it violates, and explain how you would correct it.
  \begin{center}
  \begin{tabular}{|c|c|c|}
    \hline
    \multicolumn{3}{|c|}{\textbf{Where Made}} \\\hline
    North America & Europe & Asia \\\hline
    \multicolumn{3}{|c|}{\textbf{Size}}\\\hline
    2-door & 4-door & hatchback\\\hline
  \end{tabular}
\end{center}
\boxedanswer{4in}{ %
  Some potential mistakes:
  \begin{itemize}
  \item The ``Size'' describes the number of doors, as opposed to
    'compact', 'full-size', etc. and is thus arguably incomplete. More
    importantly, the ``Size'' categories are not mutually exclusive:
    it is possible to have a hatchback car with 4-doors, for
    example. The solution is to divvy this up into proper, distinct
    size categories: [compact | mid-size | full-size | SUV/truck]
  \item The ``Where Made'' category is similarly incomplete: Brazil is
    a major car manufacturer, and Australia has been known to produce
    cars. One solution is to add additional categories, e.g.,
    [North America | Asia | Europe | South America | Other]
  \end{itemize}
}

\question
  Consider the following specification:
  \begin{minted}{python}
# Return the final price in cents after any discount:
#   - Regular: no discount
#   - Member:  10% off, but only if price is MORE THAN $100 (10000 cents)
#   - VIP:     20% off, but only if price is MORE THAN  $50  (5000 cents)
# Assume price_cents >= 0.
def calculate_discount(cust: CustomerType, price_cents: int) -> int:
    ...
  \end{minted}
  where \mintinline{python}{CustomerType} is
  \mintinline{python}{Enum("CustomerType", "Regular Member VIP")}.
\begin{parts}
\part[3] Define an input-domain model for this
  \mintinline{python}{calculate_discount} by listing a characteristic
  and its blocks for each of its parameters. Your characteristics
  should define a valid partitioning, and should leverage the expected
  behaviors described in the comment above, as opposed to relying only
  on the types of the parameters.
  \boxedanswer{3in}{ %
  \smallskip
  \begin{tabular}{|l|l|l|l|}
    \hline
    \multicolumn{4}{|l|}{\textbf{\mintinline{python}{cust}} --- customer type}\\\hline
    Regular & Member & VIP & \\\hline
    \multicolumn{4}{|l|}{\textbf{\mintinline{python}{price_cents}} --- price relative to the discount thresholds}\\\hline
    $\le 5000$ & $5000 < p \le 10000$ & $p > 10000$ & \\
    ($\le\$50$) & ($\$50$--$\$100$) & ($>\$100$) & \\\hline
  \end{tabular}
}

\part[2] Based on the price thresholds described in the comment, give
  the specific \textbf{boundary values} you would test, and state
  which fault each is meant to expose. %
  \boxedanswer{3.5in}{ %
    The thresholds are $5000$ (VIP, ``more than \$50'') and $10000$ (Member,
  ``more than \$100''). Both are strict, test on-boundary and just-above:
  \begin{itemize}
  \item \textbf{VIP threshold:} \mintinline{python}{price_cents =
      5000} (expect \emph{no} discount, since $5000 \not> 5000$) and
    \mintinline{python}{5001} (expect the 20\% discount). These are
    meant to expose an off-by-one / wrong-operator fault at the VIP
    threshold, e.g., using \mintinline{python}{>=} instead of
    \mintinline{python}{>}, which would wrongly discount at exactly
    \$50.
  \item \textbf{Member threshold:} \mintinline{python}{price_cents =
      10000} (expect \emph{no} discount) and
    \mintinline{python}{10001} (expect the 10\% discount). These
    expose the same \mintinline{python}{>=}-vs-\mintinline{python}{>}
    fault at the Member threshold.
  \item \textbf{Domain minimum:} \mintinline{python}{price_cents = 0}:
    this special value is meant to expose mishandling at the low end,
    e.g., a discount or rounding error applied to 0.
  \end{itemize}
  }

  \newpage

\part[2] A colleague partitions \mintinline{python}{cust} and
  \mintinline{python}{price_cents} independently and tests each block
  only once using the Each-Choice strategy. What assumptions does that
  testing strategy make, and how could it miss a real bug in
  \mintinline{python}{calculate_discount}? %
  \boxedanswer{4in}{ %
    Each-Choice requires only that \emph{each block of each parameter
      appear in some test}. It therefore assumes the two parameters do
    not interact, i.e., that a block of one behaves the same no matter
    which block of the other it is paired with.

    \smallskip But in \mintinline{python}{calculate_discount} these
    parameters can interact: the correct behavior at a given price
    depends on the customer type! A valid ECC set such as
  \[
    (\texttt{Regular},\ \le\$50),\quad
    (\texttt{Member},\ \$50\texttt{--}\$100),\quad
    (\texttt{VIP},\ >\$100)
  \]
  covers every block once, yet only ever tests \textbf{Member in the
    \$50--\$100 band, where the correct answer is ``no discount''
    anyway.} A bug where Member's 10\% discount never fires would pass
  all three tests, because Member is never tested above \$100, exactly
  where its discount is supposed to apply. A stronger criteria that
  pairs each customer type with the price bands around \emph{its own}
  threshold, e.g.\ Pair-Wise or Base-Choice coverage, would catch such
  an error.  }
\end{parts}

\newpage

\question Suppose we are testing a function that takes three
  parameters, which we have partitioned into the following blocks:
  \mintinline{python}{p1} $\in \{A, B\}$,\ \ \mintinline{python}{p2}
  $\in \{1, 2, 3\}$,\ \ \mintinline{python}{p3} $\in \{x, y\}$. Assume
  every combination of blocks is a valid input.
\begin{parts}
\part[3] How many tests are required to satisfy
  \textbf{All-Combinations (ACoC)}, \textbf{Each-Choice (ECC)}, and
  \textbf{Base-Choice (BCC)} coverage?  %
  \boxedanswer{2.25in}{ %
    \textbf{ACoC} $= |B_1|\cdot|B_2|\cdot|B_3| = 2\cdot 3\cdot 2 = 12$.\\
    \textbf{ECC} $= \max(2,3,2) = 3$ (one test per block of the
    largest partition,
    reusing values from the others).\\
    \textbf{BCC} $= 1 + \sum_i (|B_i|-1) = 1 + (1+2+1) = 5$: base test
    $(A,1,x)$ plus one variation per non-base block --- $(B,1,x)$,
    $(A,2,x)$, $(A,3,x)$, $(A,1,y)$.  }
\part[1] What is the minimum number of tests needed to have complete
  \textbf{Pair-Wise (PWC)} coverage for these parameters? %
  \boxedanswer{2.25in}{ %
    $6$, i.e., the product of the two largest block counts
    ($3\cdot 2$). The following set of tests achieves this:
    e.g. $(A,1,x),(A,2,y),(A,3,x),(B,1,y),(B,2,x), (B,3,y)$. Note
    that, unlike the other criteria, PWC has no closed-form solution;
    it is the size of a minimal covering test set. }
\end{parts}

\end{questions}

\end{document}
